
\newpage
%%%%%%%%%%%%%%%%%%%%%%%%%%%%%%%%%%%%%%%%%%%%%%%%%%%%%%%%%%%%%%
%%%% MA-0351 Introducción al Cálculo en Varias Variables %%%%%
%%%%%%%%%%%%%%%%%%%%%%%%%%%%%%%%%%%%%%%%%%%%%%%%%%%%%%%%%%%%%%
\subsection*{MA-0351 Introducción al Cálculo en varias variables}
\addcontentsline{toc}{subsection}{MA-0351 Introducción al Cálculo en varias variables}

\hypertarget{MA0351}{}
\begin{refsection}


\begin{table}[H]
\centering{
\begin{tabular}{|ll|}
\hline
%\multicolumn{2}{|c|}{\textbf{Nivel de virtualidad del curso:} Virtual}\\ \hline
\textbf{Año:} II  & \textbf{Requisitos:} MA-02XX Introducción al Cálculo en una variable\\ 
\textbf{Ciclo:} III  & \textbf{Correquisitos:} MA-02XX Algebra Lineal I \\ 
\textbf{Tipo de curso:} Teórico & \textbf{Horas presenciales por semana:}   \\ 
\textbf{Créditos:} 3 & \textbf{Horas de trabajo independiente por semana:}  \\ \hline
\end{tabular}
}
\end{table} 


\subsubsection*{Descripción del curso}

El Cálculo Diferencial e Integral constituye uno de los más grandes triunfos intelectuales de la historia de la humanidad. Este ha posibilitado una gran parte de los avances tecnológicos y científicos que se han tenido en los últimos siglos. Además, su descubrimiento a finales del siglo 17 por Newton y Leibniz  abrió senderos en el campo de la matemática que todavía hoy se están explorando.
Esta área del conocimiento, y los conceptos que introduce, van a aparecer repetidamente durante el resto de la carrera.

En este segundo curso de Cálculo, se presentan los conceptos fundamentales desde un punto de vista concreto, enfatizando el entendimiento por medio de la práctica. El curso empieza estudiando temas básicos sobre sucesiones y series, para luego dar inicio al estudio del cálculo diferencial e integral en varias variables, con un énfasis en el desarrollo de habilidades operacionales y del pensamiento intuitivo y geométrico con base en la exploración de una serie de conceptos y procedimientos propios del cálculo. Para lograr el desarrollo de estas habilidades e intuiciones, se priorizará el trabajo con ejemplos concretos.

Este curso se complementa con MA-2XXX Algebra Lineal I en el estudio de espacios euclídeos y sus transformaciones desde distintos puntos de vista. 



\subsubsection*{Objetivos}

\textbf{General:} Aplicar las nociones, intuiciones y conceptos básicos del cálculo diferencial e integral en varias variables para la resolución de problemas. \\

\textbf{Específicos:}

\begin{enumerate}
\item Ejemplificar casos de sucesiones, series o funciones de varias variables que cumplan propiedades predeterminadas.
\item Calcular límites de sucesiones y series a partir de propiedades básicas.
\item Distinguir convergencia de series, series de potencias y límites de funciones de varias variables.
\item Describir algebraica y geométricamente las funciones de varias variables y las regiones que definen.
\item Realizar cálculos asociados a límites y derivadas de funciones de varias variables, así como integrales múltiples, para aplicarlos a la resolución de problemas de optimización, cálculo de áreas y volúmenes, etc.
\item Aplicar la definición formal de límite en la demostración de resultados básicos asociados a sucesiones y series. 
\item Indagar en la literatura sobre temas que permitan complementar su formación.
\item Comunicar ideas matemáticas de manera clara y efectiva en forma oral y escrita.
\end{enumerate} 

\subsubsection*{Contenidos}

\begin{enumerate}
\item \textbf{Sucesiones y series (5 semanas):} 
\begin{enumerate}
    \item Definición y ejemplos de sucesiones en $\R$.
    \item Concepto intuitivo de convergencia y del límite de una sucesión convergente.
    \item Definición formal de límite de una sucesión en $\R$ con $\varepsilon$-$N$.
    \item Interpretación geométrica de la definición formal de límite.
    \item Demostración de límites concretos usando la definición formal.
    \item Propiedades básicas de límites de sucesiones (álgebra de límites, el Teorema del Encaje).
    \item Sucesiones definidas de forma recursiva.
    \item Introducción al axioma del extremo superior: Definiciones de cotas y extremos (superiores e inferiores). Cálculos con ejemplos concretos.
    \item Sucesiones monótonas, acotadas, Teorema de la Convergencia Monótona para sucesiones.
    \item Definición y ejemplos de series de números reales: serie geométrica, series telescópicas, $p$-series, 
    \item Concepto intuitivo de convergencia de series.
    \item Propiedades básicas de series (álgebra de series).
    \item Criterios de convergencia básicos: criterio de la integral, criterio de comparación, criterio de la razón, criterio para series alternantes.
    \item Convergencia absoluta y condicional.
\end{enumerate}


\item \textbf{Derivación de funciones de varias variables (5 semanas):} 
\begin{enumerate}
    \item Ejemplos básicos de funciones de varias variables $f \colon \R^n \to \R$ y sus respectivos gráficos en el caso $n = 2$. 
    \item Concepto intuitivo de límite y continuidad de una función de dos variables.
    \item Interpretación geométrica de la definición formal de límite de una función de dos variables con $\varepsilon$-$\delta$.
    \item Concepto intuitivo y geométrico de derivada parcial para funciones de varias variables.
    \item Definición y cálculo de derivadas parciales.
    \item Derivadas parciales de orden superior e igualdad de derivadas parciales mixtas (Teorema de Clairaut).
    \item Propiedades de las derivadas parciales: reglas de suma, producto, cociente y la Regla de la Cadena.
    \item Diferenciación implícita de funciones de varias variables.
    \item Concepto intuitivo y geométrico de derivadas direccionales y del vector gradiente.
    \item Definición y cálculo de derivadas direccionales. Cálculo de derivadas direccionales máximas y mínimas (direcciones de máximo y mínimo crecimiento).
    \item Cálculo de las ecuaciones de un plano tangente y una recta normal a una superficie de nivel utilizando el vector gradiente.
    \item Definición de máximos y mínimos locales y globales de una función de varias variables y sus interpretaciones geométricas en el caso de dos variables.
    \item Definición y cálculo de puntos críticos.
    \item Clasificación de puntos críticos para funciones de dos y tres variables.
    \item Definición de extremos condicionados y el Método de los Multiplicadores de Lagrange.
\end{enumerate}

\item \textbf{Integración en varias variables operacional (5 semanas):}
\begin{enumerate}
    \item Noción intuitiva y representación gráfica de una integral definida de una función de dos variables como un límite de sumas de Riemann. 
    \item Propiedades básicas de las integrales dobles (linealidad, aditividad de regiones, cotas).
    \item Cálculo de integrales dobles sobre regiones simples  por medio de integración iterada (Teorema de Fubini).
    \item Cálculo de integrales dobles en coordenadas polares.
    \item Cálculo de áreas y volúmenes usando integrales dobles.
    \item Concepto y propiedades básicas de integrales triples (linealidad, aditividad de regiones, cotas).
    \item Cálculo de integrales triples sobre regiones simples  por medio de integración iterada (Teorema de Fubini).
    \item Cálculo de integrales triples en coordenadas cilíndricas y esféricas.
    \item Cálculo de volúmenes usando integrales triples.
    \item El Teorema del Cambio de Variables en integrales múltiples.
\end{enumerate}


\end{enumerate}

\subsubsection*{Metodología}

\noindent \textbf{Lineamientos metodológicos:} La persona docente a cargo de este curso debe respetar las siguientes pautas:
\begin{enumerate}
    \item La presentación de los contenidos debe buscar un balance entre la intuición, el desarrollo de habilidades operacionales y el formalismo. Para esto se debe recurrir a ejemplos concretos con el fin de que el estudiante vaya desarrollando la madurez necesaria en el tema. 
    \item El desarrollo del axioma del extremo superior o sus equivalentes será únicamente introductorio, con el objetivo de dar las primeras nociones. El estudio del axioma del extremo superior y sus consecuencias se llevará a cabo a profundidad en el curso posterior de principios de Análisis Matemático. 
    \item Los argumentos con $\varepsilon$-$N$ se aplicarán únicamente en la demostración de las propiedades básicas de convergencia y se debe evitar entrar en la discusión de los detalles sobre propiedades conjuntistas (cardinalidad) y topológicas de $\mathbb{R}^n$.
    \item Las demostraciones que se deben realizar en este curso son las que contribuyen al desarrollo de la intuición, las habilidades operacionales y el formalismo en la persona estudiante.
    %\item Como existe codependencia entre la temática y objetivos de este curso y la del curso MA-2XXX Algebra Lineal I, debe respetarse estrictamente el orden de temas establecido en este documento.
\end{enumerate}

\textbf{Sugerencias metodológicas:} Adicionalmente, se tienen las siguientes recomendaciones para la persona docente a cargo de este curso:
\begin{enumerate}
    \item Utilizar el recurso tecnológico para reforzar distintos conceptos estudiados en el curso por medio de la visualización, cálculo y experimentación. Se pone a disposición de la persona docente un repositorio de herramientas listas para ser implementadas en el curso, tanto en las clases como para el trabajo independiente de las personas estudiantes.
    \item Aprovechar momentos adecuados del curso para motivar el estudio por medio de reseñas acerca del desarrollo histórico del cálculo y su importancia en aplicaciones dentro y fuera de la matemática
    \item En la evaluación, se sugiere utilizar estrategias que promuevan el trabajo constante de parte de las personas estudiantes a lo largo del ciclo lectivo. Por ejemplo, esto se puede hacer por medio de tareas o quizzes.
    \item Para continuar con el desarrollo de habilidades investigativas, se sugiere la implementación de pequeños proyectos de investigación. Por ejemplo, pueden ser proyectos de investigación bibliográfica en temas históricos o también proyectos cortos de aplicación del cálculo como los que se encuentran en el libro de James Stewart \cite{Ste18}.
\end{enumerate}

%\subsubsection*{Guía para las referencias}

\nocite{Abb15}
\nocite{Apo84}
\nocite{Apo85} 
\nocite{BS11}
\nocite{Apo76}
%\nocite{SRW17}
\nocite{Ste18}
\nocite{Spi08}
\nocite{Ros13}


\printbibliography[heading=subbibliography]
\end{refsection}
